\chapter{格点规范理论的表述}

1974 年，Wilson [31] 在格点上建立了欧几里得规范理论，作为研究禁闭并对 QCD 进行非微扰分析的一种手段。路径积分方法的数值实现需要以下五个步骤。
\begin{itemize}
  \item 时空的离散化。
  \item 规范与费米自由度的转写。
  \item 作用量的构造。
  \item 路径积分中积分测度的定义。
  \item 用于探测物理的算符的转写。
\end{itemize}
其中，作用量与算符的构造最为精细复杂。与其给出最新且改进过的格点作用量——那将是 Martin L\"uscher 讲义的主题——我选择遵循最初发现的道路。

\section{离散时空网格}

在 4 维欧几里得时空中有多种离散化的可能方式，其中包括超立方格点、体心立方格点 [32] 以及随机格点 [33]。其中最简单的当属各向同性超立方网格，其格距为 \(a = a_S = a_T\)，尺寸为 \(N_S \times N_S \times N_S \times N_T\)。由于该网格几乎是唯一被使用的，本讲义也仅讨论这一种。对于随机格点，即使是自由场分析也变得困难 [33]。事实上，对其它变体网格几何的探索很少，因为人们不期望能有显著的物理收益来抵消额外增加的计算复杂度。最近，涉及各向异性格点——即在空间与时间方向具有不同格距的超立方网格——的方案显示出减少离散化误差的潜力。在 18.3 节讨论胶球谱时，我将再回到这一进展。

\section{场变量 \(\psi(x)\) 与 \(A_\mu(x)\) 的格点转写}

物质场 \(\psi(x)\) 的格点转写是直接的。夸克场由定义在格点每个格位上的反对易 Grassmann 变量表示，它们属于 \(SU(3)\) 的基本表示。对一对费米变量 \(\psi\) 和 \(\bar{\psi}\) 的积分规则是 Grassmann 变量的标准规则
\begin{equation}
\begin{aligned}
\int d\psi\,d\bar{\psi} &= \int d\psi\,d\bar{\psi}\,\psi = \int d\psi\,d\bar{\psi}\,\bar{\psi} = 0 , \\
\int d\psi\,\psi &= \int d\bar{\psi}\,\bar{\psi} = i , \\
\int d\psi\,d\bar{\psi}\,\psi\,\bar{\psi} &= 1 .
\end{aligned}
\end{equation}
由于费米子作用量对 \(\psi\) 和 \(\bar{\psi}\) 都是线性的，可以用这些规则对它们积分，于是路径积分约化为仅对规范自由度的路径积分。因此，实际上人们无需担心在格点上转写 Grassmann 变量以及实现泡利不相容原理。

顺便提一下，Grassmann 变量的积分规则足以完成费米子路径积分。考虑 \(\langle \psi(x_1) \ldots\allowbreak \psi(x_n)\bar{\psi}(x_1) \ldots\allowbreak \bar{\psi}(x_n) \rangle\)。要得到非零的期望值，需要用 \(\psi\bar{\psi}\) 饱和每个格位。算符中的场所未覆盖的格位，必须来自展开式 \(\exp S_f = 1 + S_f + S_f^2/2 + \ldots\) 中的 \(S_f^{V-n}\) 项，其中 \(V\) 是格位总数。由于 \(S_f = \sum_x \bar{\psi}(x)(\not D + m)\psi(x)\) 是各格位之和，\(S_f^{V-n}\) 中 \(V-n\) 项的数目构成了一个指数级困难的组合问题。此外，对连通图有贡献的每一项都包含一条连接算符中各格位的规范链接的路径有序乘积（来自 \(S_f\) 中离散化的 \(\not D\) 项；例如见式 7.5），而费米子积分为 \(\pm m^{V-n}\)，其符号由 \(\psi\) 与 \(\bar{\psi}\) 反交换的次数决定。你应该说服自己：费米子泛函积分是对所有以最多一次触及每个格位的可能路径的求和。在枚举了所有这些路径及其权重之后，规范泛函积分现在要求对每条路径中的每个链接进行积分。这显然是一种非常困难的方法，在实践中很少被采用。

规范场的构造则不那么直观。在连续情形中，规范场 \(A_\mu(x)\) 携带 4 矢量 Lorentz 指标，并传递费米子之间的相互作用。为了转写它们，Wilson 注意到：在连续情形中，一个费米子在规范场 \(A_\mu(x)\) 存在下从格位 \(x\) 移动到 \(y\) 时，会获得由路径有序乘积给出的相位因子
\begin{equation}
\psi(y) = P\, e^{\,ig\int_x^y A_\mu(x)\,dx_\mu}\, \psi(x) .
\end{equation}
传递夸克之间相互作用的八种胶子用矩阵 \(A_\mu(x) \equiv A_\mu^a(x) \cdot \lambda^a/2\) 表示，其中群生成元 \(\lambda^a\) 按 \(\Tr \lambda^a \lambda^b = 2\delta^{ab}\) 归一化。式 6.2 提示，规范场应与连接格点上各格位的链接相联系。于是，Wilson 将路径有序乘积的离散版本与每个链接相联系
\begin{equation}
U(x, x+\hat{\mu}) \equiv U_\mu(x) = e^{iagA_\mu(x+\hat{\mu}/2)} ,
\end{equation}
其中，具体而言，平均场 \(A_\mu\) 定义在链接的中点，而 \(U\) 是一个行列式为 1 的 \(3 \times 3\) 幺正矩阵。此外，式 (2.3) 中的路径有序性规定
\begin{equation}
U(x, x-\hat{\mu}) \equiv U_{-\mu}(x) = e^{-iagA_\mu(x-\hat{\mu}/2)} = U^\dagger(x - \hat{\mu}, x) .
\end{equation}

\section{格点理论的离散对称性}

连续理论的对称群——Poincar\'e 不变性——被约化为一个离散群。在超立方格点上，只允许 \(90^\circ\) 的旋转，因此连续转动群被离散超立方群 [34] 取代。平移必须至少移动一个格点单位，因此允许的动量是离散的
\[
k = \frac{2\pi n}{La}, \qquad n = 0, 1, \ldots L
\]
或者等价地
\[
k = \pm \frac{2\pi n}{La}, \qquad n = 0, 1, \ldots L/2 .
\]
在格点上，动量按模 \(2\pi\) 守恒。

除定域规范对称性和 Poincar\'e 不变性外，格点作用量在宇称 (P)、电荷共轭 (C) 和时间反演 (T) 下也不变。场变量在这些离散变换下的行为见表 2。

\begin{table}[htbp]
  \centering
  \begin{tabular}{lccc}
    \toprule
     & \(P\) & \(C\) & \(T\) \\
    \midrule
    \(U_4(x, \tau)\) & \(U_4(-x, \tau)\) & \(U_4^{*}(x, \tau)\) & \(U_{-4}(x, -\tau)\) \\
    \(U_i(x, \tau)\) & \(U_{-i}(-x, \tau)\) & \(U_i^{*}(x, \tau)\) & \(U_i(x, -\tau)\) \\
    \(\psi(x, \tau)\) & \(\gamma_4 \psi(-x, \tau)\) & \(C\bar{\psi}^{T}(x, \tau)\) & \(\gamma_4 \gamma_5 \psi(x, -\tau)\) \\
    \(\bar{\psi}(x, \tau)\) & \(\bar{\psi}(-x, \tau)\gamma_4\) & \(-\psi^{T}(x, \tau)C^{-1}\) & \(\bar{\psi}(x, -\tau)\gamma_5 \gamma_4\) \\
    \bottomrule
  \end{tabular}
  \caption{规范与费米自由度在宇称、电荷共轭和时间反演离散变换下的行为。电荷共轭矩阵 \(C\) 满足关系 \(C\gamma_\mu C^{-1} = -\gamma_\mu^{T} = -\gamma_\mu^{*}\)，并可用 \(C = \gamma_4 \gamma_2\) 表示。注意，在欧几里得表述中 \(T\) 是线性算符。}
  \label{tab:2}
\end{table}

\section{定域规范对称性}

定域规范变换 \(V(x)\) 对变量 \(\psi(x)\) 和 \(U\) 的作用定义为
\begin{equation}
\begin{aligned}
\psi(x) &\to V(x)\psi(x) , \\
\bar{\psi}(x) &\to \bar{\psi}(x)V^\dagger(x) , \\
U_\mu(x) &\to V(x)U_\mu(x)V^\dagger(x+\hat{\mu})
\end{aligned}
\end{equation}
其中 \(V(x)\) 与 \(U\) 处于同一表示，即它是一个 \(SU(3)\) 矩阵。在这些定义下，格点上可以构造两类规范不变的对象。
\begin{itemize}
  \item 由链接的路径有序乘积构成的一条弦，其两端分别由一个费米子和一个反费米子封口，如图~\ref{fig:5}a 所示。一个简单的例子是
  \begin{equation}
  \Tr\, \bar{\psi}(x)\, U_\mu(x)\, U_\nu(x+\hat{\mu}) \ldots U_\rho(y-\hat{\rho})\, \psi(y)
  \end{equation}
  其中迹是对颜色指标求迹。我将用 \emph{弦}（string）一词指称其推广形式：如式 6.6 所定义的单个自旋与颜色迹，或若干条路径有序且无空间间断的基本弦的乘积。在具有周期性边界条件的格点上，如果弦横跨整个格点并由周期性闭合，则无需 \(\psi\)、\(\bar{\psi}\) 封口。这样的弦被称为威尔逊线/波利亚科夫线。
  \item 封闭的威尔逊圈，如图~\ref{fig:5}b 所示。最简单的例子是方格，即一个 \(1 \times 1\) 圈，
  \begin{equation}
  W^{\mathrm{1\times 1}}_{\mu\nu} = \Re\, \Tr\, U_\mu(x)\, U_\nu(x+\hat{\mu})\, U_\mu^\dagger(x+\hat{\nu})\, U_\nu^\dagger(x) .
  \end{equation}
\end{itemize}
除非另有说明，每当我使用威尔逊圈一词时，均假定已取迹的实部 Re Tr。对于 \(SU(N \ge 3)\)，基本表示中任何威尔逊圈的迹都是复数，两种可能的路径有序给出互为复共轭的值。因此，取迹保证了规范不变性，而取实部等价于对该圈及其电荷共轭取平均。

\textit{练习：} 请自行确认，不存在任何其它独立的规范不变量。

\begin{figure}[htbp]
  \centering
  \includegraphics[width=0.8\textwidth]{images/fig5.png}
  \caption{两类规范不变量。a) 由一系列有序的 $U$ 构成、两端分别以费米子和反费米子封口的弦；b) 封闭的威尔逊圈。}
  \label{fig:5}
\end{figure}

\textit{练习：} 请自行确认，为了保持规范不变性，格点理论必须用链接变量 \(U_\mu\) 来表述，而不是直接用 \(A_\mu\)。

因此，规范不变的作用量必须由圈和弦构造而成。不过，请注意这两类基本的规范不变对象可以有任意的大小和形状。此外，它们可以处于颜色 \(SU(3)\) 的任何表示中。唯一必要的约束是，它必须在 \(a \to 0\) 极限下给出我们所熟悉的连续理论。希望在任意 \(a\) 下都保持规范不变性的原因是：否则将会有多得多的参数需要调节（零胶子质量，以及夸克-胶子、三胶子和四胶子耦合的相等性都是规范对称性的推论），并且在 \(a\) 的任意给定阶都会出现多得多的算符（这些是无关算符，但在有限 \(a\) 下会产生离散化误差，如第 8 节所述）。事实上，人们希望在 \(a\) 的一切取值下都尽可能多地保持理论的对称性，而不仅仅是在连续极限下。结果表明，保持定域规范不变性相对简单；然而，保持手征对称性——它在轻夸克唯象学中已被证明如此重要——以及保持 Dirac 费米子与格点费米子自由度之间的一一对应关系，却一直未获成功。这一点将在第 9 节讨论。此外，离散化事实上破坏了连续的转动、Lorentz 和平移对称性。改进离散化方案的关键思想——将在第 8、12 和 13 节讨论——是以携带足够多 $a$ 幂次的算符来实现这些对称性的破缺，从而使它们的影响非常局域，因而对物理（长距离）量很小。
